\par\Needspace{6\baselineskip}\medskip
\begin{center}\hypertarget{native-173-article126}{}\pdfbookmark[2]{126.}{article-126}126.\end{center}
Methodus praecedens ad \emph{tertiam} illico sternit viam. Determinentur, perinde
ut ante, e distantiis corporis coelestis a terra in observationibus extremis
longitudines respondentes in orbita cum radiis vectoribus. Adiumento positionis
plani orbitae, quam hic calculus suppeditaverit, eruatur ex observatione
media longitudo in orbita atque radius vector. Tunc autem computentur ex
his tribus locis heliocentricis elementa reliqua per problema in artt. 82, 83
tractatum, quae operatio ab observationum temporibus independens erit. Hoc
itaque modo innotescent tres anomaliae mediae atque motus diurnus, unde
ipsa temporum intervalla inter observationem primam et secundam atque inter
secundam et tertiam computare licebit. Horum differentiae ab intervallis veris
pro $X$ et $Y$ accipientur.

Haec methodus minus idonea esset, quoties motus heliocentricus arcum
exiguum tantum complectitur. In tali enim casu ista orbitae determinatio (ut
iam in art. 82 monuimus) a quantitatibus tertii ordinis pendet, adeoque praecisionem
sufficientem non admittit. Variationes levissimae in valoribus ipsarum
$x,y$ producere possent variationes permagnas in elementis adeoque etiam
in valoribus ipsarum $X,Y$ neque has illis proportionales supponere liceret.
Quoties autem tres loci motum heliocentricum considerabilem subtendunt, methodi
usus utique succedet optime, siquidem exceptionibus in artt. praec. \SourceLeaf{174}explicatis
haud turbetur, ad quas manifesto in hac quoque methodo respiciendum
erit.
